\chapter{Motivação e objetivos}\label{chap:problema}
    Texto introdutório de seção/capítulo. Deve-se descrever em linhas gerais o que será abordado nesta seção/capítulo.

\section{Motivação}
    Na motivação, deve-se escrever sobre o porquê que determinado trabalho está sendo projetado. Num trabalho final, deve-se apresentar a solução para um problema \citep{WAZLAWICK2020}.  Deve-se definir da forma mais clara possível o problema a ser solucionado, indicando \textbf{onde} e \textbf{de que forma} ele impacta em sua área específica. A revisão de literatura (\cref{chap:revisao}) pode (deve) auxiliar na escrita desta seção. 
    
\section{Objetivo Geral}\label{sec:obj_geral}

O Objetivo Geral é o mais amplo de seu trabalho. Considere que se tudo proposto foi executado corretamente e deu certo, você consegue alcançar seu objetivo. 

\section{Objetivos Específicos}\label{sec:obj_espec}
Nos objetivos específicos pode-se relacionar todos os pequenos objetivos a serem realizados, para no final conseguir alcançar o objetivo geral. Os objetivos (geral e específicos) devem ser expressos cujos sucesso possa ser provado, de forma não trivial \cite{WAZLAWICK2020}. Uma boa sugestão é iniciar os objetivos com verbos.
\begin{itemize}[leftmargin=5em]
    \item Estudar sobre o tema de pesquisa;
    \item Definir modelo de solução;
    \item Desenvolver a solução com uso de ....
    \item Validar ....
    \item objetivo específico ...
\end{itemize}    

\section{Hipóteses [Opcional]}\label{sec:hipoteses}
Caso seu trabalho contenha hipóteses, descreva as mesmas aqui.
    


